% =======================================================================
% preamble.tex -- Metropolis-themed Beamer preamble for technical paper talks
% Intended to be included from <paper-slug>-Slides.tex via % =======================================================================
% preamble.tex -- Metropolis-themed Beamer preamble for technical paper talks
% Intended to be included from <paper-slug>-Slides.tex via % =======================================================================
% preamble.tex -- Metropolis-themed Beamer preamble for technical paper talks
% Intended to be included from <paper-slug>-Slides.tex via % =======================================================================
% preamble.tex -- Metropolis-themed Beamer preamble for technical paper talks
% Intended to be included from <paper-slug>-Slides.tex via \input{preamble.tex}.
% Compile with pdflatex by default; xelatex/lualatex also work.
% =======================================================================

% ---- Document class is set in <paper-slug>-Slides.tex: \documentclass[aspectratio=169]{beamer}

% ---- Engine -----------------------------------------------------------
\usepackage{iftex}

% ---- Theme ------------------------------------------------------------
\usetheme{metropolis}
\setbeamercovered{transparent}
\metroset{
  progressbar=frametitle,
  numbering=fraction,
  block=fill,
  sectionpage=progressbar,
  subsectionpage=progressbar
}
\usefonttheme{professionalfonts}
\ifPDFTeX
  \usepackage[T1]{fontenc}
  \usepackage[utf8]{inputenc}
  \usepackage[sfdefault,light]{FiraSans}
\else
  \setsansfont{Fira Sans Light}[
    Ligatures=TeX,
    ItalicFont={Fira Sans Light Italic},
    BoldFont={Fira Sans},
    BoldItalicFont={Fira Sans Italic}
  ]
\fi

% ---- Math / symbols ---------------------------------------------------
\usepackage{amsmath,amssymb,amsthm,mathtools}
\usepackage{bm}
\usepackage{bbm}
\usepackage{dsfont}

% ---- Algorithms -------------------------------------------------------
\usepackage[ruled,linesnumbered,vlined]{algorithm2e}
% \usepackage{algpseudocode}  % alternative

% ---- Graphics / tables ------------------------------------------------
\usepackage{graphicx}
\usepackage{booktabs}
\usepackage{multirow}
\usepackage{makecell}
\graphicspath{{figs/}}

% ---- TikZ / PGF -------------------------------------------------------
\usepackage{tikz}
\usetikzlibrary{positioning,arrows.meta,calc,shapes.geometric,fit,backgrounds}
\usepackage{pgfplots}
% Metropolis loads its Tol pgfplots theme late and resets compat to 1.9.
\AtEndPreamble{\pgfplotsset{compat=1.18}}

% ---- References -------------------------------------------------------
\usepackage[backend=biber,style=numeric,sorting=none,maxbibnames=99]{biblatex}
\addbibresource{references.bib}

% ---- Theorem environments --------------------------------------------
% Keep paper numbering; if authors' Thm 3.2 appears, we cite as "(Thm 3.2)".
\makeatletter
\theoremstyle{definition}
\@ifundefined{assumption}{\newtheorem{assumption}{Assumption}}{}
\@ifundefined{definition}{\newtheorem{definition}{Definition}[section]}{}

\theoremstyle{plain}
\@ifundefined{theorem}{\newtheorem{theorem}{Theorem}[section]}{}
\@ifundefined{lemma}{\newtheorem{lemma}[theorem]{Lemma}}{}
\@ifundefined{proposition}{\newtheorem{proposition}[theorem]{Proposition}}{}
\@ifundefined{corollary}{\newtheorem{corollary}[theorem]{Corollary}}{}

\theoremstyle{remark}
\@ifundefined{remark}{\newtheorem*{remark}{Remark}}{}
\@ifundefined{question}{\newtheorem*{question}{Question}}{}       % for socratic questions
\makeatother

% ---- Custom block colors (second accent for definitions vs theorems) --
\setbeamercolor{block title}{bg=mDarkTeal,fg=white}
\setbeamercolor{block body}{bg=mDarkTeal!10,fg=black}

% Add modest outer and inner padding to Metropolis callout blocks.
\makeatletter
\newlength{\BY@blockmargin}
\setlength{\BY@blockmargin}{0.45em}
\setlength{\metropolis@blocksep}{0.95ex}
\setlength{\metropolis@blockadjust}{0.30ex}

\newcommand{\BY@blockbegin}[1]{%
  \par\noindent\hfill%
  \begin{minipage}{\dimexpr\linewidth-2\BY@blockmargin\relax}%
  \metropolis@block{#1}%
}
\newcommand{\BY@blockend}{%
  \end{beamercolorbox}\vspace*{0.2ex}%
  \end{minipage}\hfill\par%
}
\setbeamertemplate{block begin}{\BY@blockbegin{}}
\setbeamertemplate{block alerted begin}{\BY@blockbegin{ alerted}}
\setbeamertemplate{block example begin}{\BY@blockbegin{ example}}
\setbeamertemplate{block end}{\BY@blockend}
\setbeamertemplate{block alerted end}{\BY@blockend}
\setbeamertemplate{block example end}{\BY@blockend}
\makeatother

% ---- Common math macros ----------------------------------------------
\newcommand{\RR}{\mathbb{R}}
\newcommand{\NN}{\mathbb{N}}
\newcommand{\ZZ}{\mathbb{Z}}
\newcommand{\EE}{\mathbb{E}}
\newcommand{\PP}{\mathbb{P}}
\newcommand{\Var}{\mathrm{Var}}
\newcommand{\Cov}{\mathrm{Cov}}
\newcommand{\one}{\mathbf{1}}

\newcommand{\eps}{\varepsilon}
\newcommand{\lam}{\lambda}
\newcommand{\grad}{\nabla}
\newcommand{\inner}[2]{\langle #1,\, #2 \rangle}
\newcommand{\norm}[1]{\left\|#1\right\|}
\newcommand{\abs}[1]{\left|#1\right|}

\DeclareMathOperator*{\argmin}{arg\,min}
\DeclareMathOperator*{\argmax}{arg\,max}
\DeclareMathOperator{\prox}{prox}
\DeclareMathOperator{\dom}{dom}
\DeclareMathOperator{\tr}{tr}
\DeclareMathOperator{\rank}{rank}
\DeclareMathOperator{\diag}{diag}

% ---- Discussion cue environment --------------------------------------
% Use \begin{discuss}{title}...\end{discuss} for group-meeting hooks.
\newenvironment{discuss}[1]{%
  \begin{alertblock}{\textbf{Discussion:} #1}%
}{%
  \end{alertblock}%
}

% ---- Minor spacing fixes ---------------------------------------------
\setbeamertemplate{itemize items}[circle]
\setlength{\leftmargini}{1em}
.
% Compile with pdflatex by default; xelatex/lualatex also work.
% =======================================================================

% ---- Document class is set in <paper-slug>-Slides.tex: \documentclass[aspectratio=169]{beamer}

% ---- Engine -----------------------------------------------------------
\usepackage{iftex}

% ---- Theme ------------------------------------------------------------
\usetheme{metropolis}
\setbeamercovered{transparent}
\metroset{
  progressbar=frametitle,
  numbering=fraction,
  block=fill,
  sectionpage=progressbar,
  subsectionpage=progressbar
}
\usefonttheme{professionalfonts}
\ifPDFTeX
  \usepackage[T1]{fontenc}
  \usepackage[utf8]{inputenc}
  \usepackage[sfdefault,light]{FiraSans}
\else
  \setsansfont{Fira Sans Light}[
    Ligatures=TeX,
    ItalicFont={Fira Sans Light Italic},
    BoldFont={Fira Sans},
    BoldItalicFont={Fira Sans Italic}
  ]
\fi

% ---- Math / symbols ---------------------------------------------------
\usepackage{amsmath,amssymb,amsthm,mathtools}
\usepackage{bm}
\usepackage{bbm}
\usepackage{dsfont}

% ---- Algorithms -------------------------------------------------------
\usepackage[ruled,linesnumbered,vlined]{algorithm2e}
% \usepackage{algpseudocode}  % alternative

% ---- Graphics / tables ------------------------------------------------
\usepackage{graphicx}
\usepackage{booktabs}
\usepackage{multirow}
\usepackage{makecell}
\graphicspath{{figs/}}

% ---- TikZ / PGF -------------------------------------------------------
\usepackage{tikz}
\usetikzlibrary{positioning,arrows.meta,calc,shapes.geometric,fit,backgrounds}
\usepackage{pgfplots}
% Metropolis loads its Tol pgfplots theme late and resets compat to 1.9.
\AtEndPreamble{\pgfplotsset{compat=1.18}}

% ---- References -------------------------------------------------------
\usepackage[backend=biber,style=numeric,sorting=none,maxbibnames=99]{biblatex}
\addbibresource{references.bib}

% ---- Theorem environments --------------------------------------------
% Keep paper numbering; if authors' Thm 3.2 appears, we cite as "(Thm 3.2)".
\makeatletter
\theoremstyle{definition}
\@ifundefined{assumption}{\newtheorem{assumption}{Assumption}}{}
\@ifundefined{definition}{\newtheorem{definition}{Definition}[section]}{}

\theoremstyle{plain}
\@ifundefined{theorem}{\newtheorem{theorem}{Theorem}[section]}{}
\@ifundefined{lemma}{\newtheorem{lemma}[theorem]{Lemma}}{}
\@ifundefined{proposition}{\newtheorem{proposition}[theorem]{Proposition}}{}
\@ifundefined{corollary}{\newtheorem{corollary}[theorem]{Corollary}}{}

\theoremstyle{remark}
\@ifundefined{remark}{\newtheorem*{remark}{Remark}}{}
\@ifundefined{question}{\newtheorem*{question}{Question}}{}       % for socratic questions
\makeatother

% ---- Custom block colors (second accent for definitions vs theorems) --
\setbeamercolor{block title}{bg=mDarkTeal,fg=white}
\setbeamercolor{block body}{bg=mDarkTeal!10,fg=black}

% Add modest outer and inner padding to Metropolis callout blocks.
\makeatletter
\newlength{\BY@blockmargin}
\setlength{\BY@blockmargin}{0.45em}
\setlength{\metropolis@blocksep}{0.95ex}
\setlength{\metropolis@blockadjust}{0.30ex}

\newcommand{\BY@blockbegin}[1]{%
  \par\noindent\hfill%
  \begin{minipage}{\dimexpr\linewidth-2\BY@blockmargin\relax}%
  \metropolis@block{#1}%
}
\newcommand{\BY@blockend}{%
  \end{beamercolorbox}\vspace*{0.2ex}%
  \end{minipage}\hfill\par%
}
\setbeamertemplate{block begin}{\BY@blockbegin{}}
\setbeamertemplate{block alerted begin}{\BY@blockbegin{ alerted}}
\setbeamertemplate{block example begin}{\BY@blockbegin{ example}}
\setbeamertemplate{block end}{\BY@blockend}
\setbeamertemplate{block alerted end}{\BY@blockend}
\setbeamertemplate{block example end}{\BY@blockend}
\makeatother

% ---- Common math macros ----------------------------------------------
\newcommand{\RR}{\mathbb{R}}
\newcommand{\NN}{\mathbb{N}}
\newcommand{\ZZ}{\mathbb{Z}}
\newcommand{\EE}{\mathbb{E}}
\newcommand{\PP}{\mathbb{P}}
\newcommand{\Var}{\mathrm{Var}}
\newcommand{\Cov}{\mathrm{Cov}}
\newcommand{\one}{\mathbf{1}}

\newcommand{\eps}{\varepsilon}
\newcommand{\lam}{\lambda}
\newcommand{\grad}{\nabla}
\newcommand{\inner}[2]{\langle #1,\, #2 \rangle}
\newcommand{\norm}[1]{\left\|#1\right\|}
\newcommand{\abs}[1]{\left|#1\right|}

\DeclareMathOperator*{\argmin}{arg\,min}
\DeclareMathOperator*{\argmax}{arg\,max}
\DeclareMathOperator{\prox}{prox}
\DeclareMathOperator{\dom}{dom}
\DeclareMathOperator{\tr}{tr}
\DeclareMathOperator{\rank}{rank}
\DeclareMathOperator{\diag}{diag}

% ---- Discussion cue environment --------------------------------------
% Use \begin{discuss}{title}...\end{discuss} for group-meeting hooks.
\newenvironment{discuss}[1]{%
  \begin{alertblock}{\textbf{Discussion:} #1}%
}{%
  \end{alertblock}%
}

% ---- Minor spacing fixes ---------------------------------------------
\setbeamertemplate{itemize items}[circle]
\setlength{\leftmargini}{1em}
.
% Compile with pdflatex by default; xelatex/lualatex also work.
% =======================================================================

% ---- Document class is set in <paper-slug>-Slides.tex: \documentclass[aspectratio=169]{beamer}

% ---- Engine -----------------------------------------------------------
\usepackage{iftex}

% ---- Theme ------------------------------------------------------------
\usetheme{metropolis}
\setbeamercovered{transparent}
\metroset{
  progressbar=frametitle,
  numbering=fraction,
  block=fill,
  sectionpage=progressbar,
  subsectionpage=progressbar
}
\usefonttheme{professionalfonts}
\ifPDFTeX
  \usepackage[T1]{fontenc}
  \usepackage[utf8]{inputenc}
  \usepackage[sfdefault,light]{FiraSans}
\else
  \setsansfont{Fira Sans Light}[
    Ligatures=TeX,
    ItalicFont={Fira Sans Light Italic},
    BoldFont={Fira Sans},
    BoldItalicFont={Fira Sans Italic}
  ]
\fi

% ---- Math / symbols ---------------------------------------------------
\usepackage{amsmath,amssymb,amsthm,mathtools}
\usepackage{bm}
\usepackage{bbm}
\usepackage{dsfont}

% ---- Algorithms -------------------------------------------------------
\usepackage[ruled,linesnumbered,vlined]{algorithm2e}
% \usepackage{algpseudocode}  % alternative

% ---- Graphics / tables ------------------------------------------------
\usepackage{graphicx}
\usepackage{booktabs}
\usepackage{multirow}
\usepackage{makecell}
\graphicspath{{figs/}}

% ---- TikZ / PGF -------------------------------------------------------
\usepackage{tikz}
\usetikzlibrary{positioning,arrows.meta,calc,shapes.geometric,fit,backgrounds}
\usepackage{pgfplots}
% Metropolis loads its Tol pgfplots theme late and resets compat to 1.9.
\AtEndPreamble{\pgfplotsset{compat=1.18}}

% ---- References -------------------------------------------------------
\usepackage[backend=biber,style=numeric,sorting=none,maxbibnames=99]{biblatex}
\addbibresource{references.bib}

% ---- Theorem environments --------------------------------------------
% Keep paper numbering; if authors' Thm 3.2 appears, we cite as "(Thm 3.2)".
\makeatletter
\theoremstyle{definition}
\@ifundefined{assumption}{\newtheorem{assumption}{Assumption}}{}
\@ifundefined{definition}{\newtheorem{definition}{Definition}[section]}{}

\theoremstyle{plain}
\@ifundefined{theorem}{\newtheorem{theorem}{Theorem}[section]}{}
\@ifundefined{lemma}{\newtheorem{lemma}[theorem]{Lemma}}{}
\@ifundefined{proposition}{\newtheorem{proposition}[theorem]{Proposition}}{}
\@ifundefined{corollary}{\newtheorem{corollary}[theorem]{Corollary}}{}

\theoremstyle{remark}
\@ifundefined{remark}{\newtheorem*{remark}{Remark}}{}
\@ifundefined{question}{\newtheorem*{question}{Question}}{}       % for socratic questions
\makeatother

% ---- Custom block colors (second accent for definitions vs theorems) --
\setbeamercolor{block title}{bg=mDarkTeal,fg=white}
\setbeamercolor{block body}{bg=mDarkTeal!10,fg=black}

% Add modest outer and inner padding to Metropolis callout blocks.
\makeatletter
\newlength{\BY@blockmargin}
\setlength{\BY@blockmargin}{0.45em}
\setlength{\metropolis@blocksep}{0.95ex}
\setlength{\metropolis@blockadjust}{0.30ex}

\newcommand{\BY@blockbegin}[1]{%
  \par\noindent\hfill%
  \begin{minipage}{\dimexpr\linewidth-2\BY@blockmargin\relax}%
  \metropolis@block{#1}%
}
\newcommand{\BY@blockend}{%
  \end{beamercolorbox}\vspace*{0.2ex}%
  \end{minipage}\hfill\par%
}
\setbeamertemplate{block begin}{\BY@blockbegin{}}
\setbeamertemplate{block alerted begin}{\BY@blockbegin{ alerted}}
\setbeamertemplate{block example begin}{\BY@blockbegin{ example}}
\setbeamertemplate{block end}{\BY@blockend}
\setbeamertemplate{block alerted end}{\BY@blockend}
\setbeamertemplate{block example end}{\BY@blockend}
\makeatother

% ---- Common math macros ----------------------------------------------
\newcommand{\RR}{\mathbb{R}}
\newcommand{\NN}{\mathbb{N}}
\newcommand{\ZZ}{\mathbb{Z}}
\newcommand{\EE}{\mathbb{E}}
\newcommand{\PP}{\mathbb{P}}
\newcommand{\Var}{\mathrm{Var}}
\newcommand{\Cov}{\mathrm{Cov}}
\newcommand{\one}{\mathbf{1}}

\newcommand{\eps}{\varepsilon}
\newcommand{\lam}{\lambda}
\newcommand{\grad}{\nabla}
\newcommand{\inner}[2]{\langle #1,\, #2 \rangle}
\newcommand{\norm}[1]{\left\|#1\right\|}
\newcommand{\abs}[1]{\left|#1\right|}

\DeclareMathOperator*{\argmin}{arg\,min}
\DeclareMathOperator*{\argmax}{arg\,max}
\DeclareMathOperator{\prox}{prox}
\DeclareMathOperator{\dom}{dom}
\DeclareMathOperator{\tr}{tr}
\DeclareMathOperator{\rank}{rank}
\DeclareMathOperator{\diag}{diag}

% ---- Discussion cue environment --------------------------------------
% Use \begin{discuss}{title}...\end{discuss} for group-meeting hooks.
\newenvironment{discuss}[1]{%
  \begin{alertblock}{\textbf{Discussion:} #1}%
}{%
  \end{alertblock}%
}

% ---- Minor spacing fixes ---------------------------------------------
\setbeamertemplate{itemize items}[circle]
\setlength{\leftmargini}{1em}
.
% Compile with pdflatex by default; xelatex/lualatex also work.
% =======================================================================

% ---- Document class is set in <paper-slug>-Slides.tex: \documentclass[aspectratio=169]{beamer}

% ---- Engine -----------------------------------------------------------
\usepackage{iftex}

% ---- Theme ------------------------------------------------------------
\usetheme{metropolis}
\setbeamercovered{transparent}
\metroset{
  progressbar=frametitle,
  numbering=fraction,
  block=fill,
  sectionpage=progressbar,
  subsectionpage=progressbar
}
\usefonttheme{professionalfonts}
\ifPDFTeX
  \usepackage[T1]{fontenc}
  \usepackage[utf8]{inputenc}
  \usepackage[sfdefault,light]{FiraSans}
\else
  \setsansfont{Fira Sans Light}[
    Ligatures=TeX,
    ItalicFont={Fira Sans Light Italic},
    BoldFont={Fira Sans},
    BoldItalicFont={Fira Sans Italic}
  ]
\fi

% ---- Math / symbols ---------------------------------------------------
\usepackage{amsmath,amssymb,amsthm,mathtools}
\usepackage{bm}
\usepackage{bbm}
\usepackage{dsfont}

% ---- Algorithms -------------------------------------------------------
\usepackage[ruled,linesnumbered,vlined]{algorithm2e}
% \usepackage{algpseudocode}  % alternative

% ---- Graphics / tables ------------------------------------------------
\usepackage{graphicx}
\usepackage{booktabs}
\usepackage{multirow}
\usepackage{makecell}
\graphicspath{{figs/}}

% ---- TikZ / PGF -------------------------------------------------------
\usepackage{tikz}
\usetikzlibrary{positioning,arrows.meta,calc,shapes.geometric,fit,backgrounds}
\usepackage{pgfplots}
% Metropolis loads its Tol pgfplots theme late and resets compat to 1.9.
\AtEndPreamble{\pgfplotsset{compat=1.18}}

% ---- References -------------------------------------------------------
\usepackage[backend=biber,style=numeric,sorting=none,maxbibnames=99]{biblatex}
\addbibresource{references.bib}

% ---- Theorem environments --------------------------------------------
% Keep paper numbering; if authors' Thm 3.2 appears, we cite as "(Thm 3.2)".
\makeatletter
\theoremstyle{definition}
\@ifundefined{assumption}{\newtheorem{assumption}{Assumption}}{}
\@ifundefined{definition}{\newtheorem{definition}{Definition}[section]}{}

\theoremstyle{plain}
\@ifundefined{theorem}{\newtheorem{theorem}{Theorem}[section]}{}
\@ifundefined{lemma}{\newtheorem{lemma}[theorem]{Lemma}}{}
\@ifundefined{proposition}{\newtheorem{proposition}[theorem]{Proposition}}{}
\@ifundefined{corollary}{\newtheorem{corollary}[theorem]{Corollary}}{}

\theoremstyle{remark}
\@ifundefined{remark}{\newtheorem*{remark}{Remark}}{}
\@ifundefined{question}{\newtheorem*{question}{Question}}{}       % for socratic questions
\makeatother

% ---- Custom block colors (second accent for definitions vs theorems) --
\setbeamercolor{block title}{bg=mDarkTeal,fg=white}
\setbeamercolor{block body}{bg=mDarkTeal!10,fg=black}

% Add modest outer and inner padding to Metropolis callout blocks.
\makeatletter
\newlength{\BY@blockmargin}
\setlength{\BY@blockmargin}{0.45em}
\setlength{\metropolis@blocksep}{0.95ex}
\setlength{\metropolis@blockadjust}{0.30ex}

\newcommand{\BY@blockbegin}[1]{%
  \par\noindent\hfill%
  \begin{minipage}{\dimexpr\linewidth-2\BY@blockmargin\relax}%
  \metropolis@block{#1}%
}
\newcommand{\BY@blockend}{%
  \end{beamercolorbox}\vspace*{0.2ex}%
  \end{minipage}\hfill\par%
}
\setbeamertemplate{block begin}{\BY@blockbegin{}}
\setbeamertemplate{block alerted begin}{\BY@blockbegin{ alerted}}
\setbeamertemplate{block example begin}{\BY@blockbegin{ example}}
\setbeamertemplate{block end}{\BY@blockend}
\setbeamertemplate{block alerted end}{\BY@blockend}
\setbeamertemplate{block example end}{\BY@blockend}
\makeatother

% ---- Common math macros ----------------------------------------------
\newcommand{\RR}{\mathbb{R}}
\newcommand{\NN}{\mathbb{N}}
\newcommand{\ZZ}{\mathbb{Z}}
\newcommand{\EE}{\mathbb{E}}
\newcommand{\PP}{\mathbb{P}}
\newcommand{\Var}{\mathrm{Var}}
\newcommand{\Cov}{\mathrm{Cov}}
\newcommand{\one}{\mathbf{1}}

\newcommand{\eps}{\varepsilon}
\newcommand{\lam}{\lambda}
\newcommand{\grad}{\nabla}
\newcommand{\inner}[2]{\langle #1,\, #2 \rangle}
\newcommand{\norm}[1]{\left\|#1\right\|}
\newcommand{\abs}[1]{\left|#1\right|}

\DeclareMathOperator*{\argmin}{arg\,min}
\DeclareMathOperator*{\argmax}{arg\,max}
\DeclareMathOperator{\prox}{prox}
\DeclareMathOperator{\dom}{dom}
\DeclareMathOperator{\tr}{tr}
\DeclareMathOperator{\rank}{rank}
\DeclareMathOperator{\diag}{diag}

% ---- Discussion cue environment --------------------------------------
% Use \begin{discuss}{title}...\end{discuss} for group-meeting hooks.
\newenvironment{discuss}[1]{%
  \begin{alertblock}{\textbf{Discussion:} #1}%
}{%
  \end{alertblock}%
}

% ---- Minor spacing fixes ---------------------------------------------
\setbeamertemplate{itemize items}[circle]
\setlength{\leftmargini}{1em}
